\chapter{Chapter06 — Complete Quadratic Decomposition}
%%% generated by the blueprint pipeline from FABL.Chapter06.F₂Polynomials.CompleteQuadraticDecomposition
\section{Overview}

Declarations of the Lean module \texttt{FABL.Chapter06.F₂Polynomials.CompleteQuadraticDecomposition}.

\section{Declarations}

\begin{definition}[Splitting a cube into two blocks]
\lean{FABL.cubeAppendLinearEquiv}
\label{FABL.cubeAppendLinearEquiv}
\leanok
\uses{FABL.F₂Cube, FABL.𝔽₂}
For natural numbers $m$ and $p$, the vector space of functions from an $(m+p)$-element
index set to $\mathbb{F}_2$ is linearly equivalent to the product of the corresponding
spaces of functions on $m$ and $p$ coordinates, by splitting the coordinates into two
consecutive blocks.
\end{definition}

\begin{definition}[Splitting off the last two coordinates]
\lean{FABL.splitLastTwoLinearEquiv}
\label{FABL.splitLastTwoLinearEquiv}
\leanok
\uses{FABL.F₂Cube, FABL.cubeAppendLinearEquiv, FABL.𝔽₂}
For a natural number $m$, the space of $\mathbb{F}_2$-valued vectors with $m+2$
coordinates is linearly equivalent to the product of the space with $m$ coordinates and
two copies of $\mathbb{F}_2$, by splitting off the last two coordinates.
\end{definition}

\begin{definition}[Appending two coordinates]
\lean{FABL.appendTwo}
\label{FABL.appendTwo}
\leanok
\uses{FABL.F₂Cube, FABL.𝔽₂}
For a vector $x$ of $m$ coordinates over $\mathbb{F}_2$ and scalars
$a,b\in\mathbb{F}_2$, append $a$ and $b$ to $x$ to obtain a vector of $m+2$ coordinates.
\end{definition}

\begin{theorem}[Splitting Off the Last Two Coordinates]
\lean{FABL.splitLastTwoLinearEquiv_appendTwo}
\label{FABL.splitLastTwoLinearEquiv_appendTwo}
\leanok
\uses{FABL.F₂Cube, FABL.appendTwo, FABL.cubeAppendLinearEquiv, FABL.splitLastTwoLinearEquiv, FABL.𝔽₂}
\difficulty{1}
Let $m$ be a natural number, let $x\in\bbf_2^m$, and let $a,b\in\bbf_2$. Under the
linear coordinate identification of $\bbf_2^{m+2}$ with
$\bbf_2^m\times(\bbf_2\times\bbf_2)$, appending $a,b$ to $x$ maps to $(x,(a,b))$.
\end{theorem}

\begin{theorem}[Reconstructing a Vector from Its Coordinates]
\lean{FABL.splitLastTwoLinearEquiv_symm_apply}
\label{FABL.splitLastTwoLinearEquiv_symm_apply}
\leanok
\uses{FABL.F₂Cube, FABL.appendTwo, FABL.splitLastTwoLinearEquiv, FABL.𝔽₂}
\difficulty{3}
Let $m$ be a natural number, $x\in\bbf_2^m$, and $a,b\in\bbf_2$. The inverse of the
coordinate identification $\bbf_2^{m+2}\simeq\bbf_2^m\times(\bbf_2\times\bbf_2)$ sends
$(x,(a,b))$ to the vector obtained by appending $a,b$ to $x$.
\end{theorem}

\begin{definition}[Sum of coordinates]
\lean{FABL.coordinateSumLinearMap}
\label{FABL.coordinateSumLinearMap}
\leanok
\uses{FABL.F₂Cube, FABL.𝔽₂}
For a natural number $m$, the coordinate-sum linear functional on the space of
$\mathbb{F}_2$-valued vectors with $m$ coordinates sends a vector $x$ to
$\sum_{i=1}^{m}x_i$.
\end{definition}

\begin{theorem}[Coordinate sum linear map evaluation]
\lean{FABL.coordinateSumLinearMap_apply}
\label{FABL.coordinateSumLinearMap_apply}
\leanok
\uses{FABL.F₂Cube, FABL.coordinateSumLinearMap}
\difficulty{1}
For any natural number $m$ and any function $x : \\{0,1,\\ldots,m-1\\} \\to
\\mathbb{F}_2$, the coordinate-sum linear map sends $x$ to $\\sum_{i=0}^{m-1} x(i)$.
\end{theorem}

\begin{definition}[Shearing the last two coordinates]
\lean{FABL.tailShearLinearMap}
\label{FABL.tailShearLinearMap}
\leanok
\uses{FABL.F₂Cube, FABL.coordinateSumLinearMap, FABL.𝔽₂}
For a natural number $m$, define a linear map on the product of the $m$-coordinate space
with $\mathbb{F}_2^2$ by sending $(x,(a,b))$ to $(x,(a+s(x),b+s(x)))$, where $s(x)$ is
the sum of the coordinates of $x$.
\end{definition}

\begin{theorem}[Involutivity of the tail shear]
\lean{FABL.tailShearLinearMap_involutive}
\label{FABL.tailShearLinearMap_involutive}
\leanok
\uses{FABL.tailShearLinearMap, FABL.𝔽₂}
\difficulty{4}
For every natural number $m$, applying the tail shear linear map twice to any element of
$(\\mathbb{F}_2)^m\\times(\\mathbb{F}_2\\times\\mathbb{F}_2)$ returns that element.
\end{theorem}

\begin{definition}[Shear automorphism]
\lean{FABL.tailShearLinearEquiv}
\label{FABL.tailShearLinearEquiv}
\leanok
\uses{FABL.F₂Cube, FABL.tailShearLinearMap, FABL.tailShearLinearMap_involutive, FABL.𝔽₂}
For a natural number $m$, the map $(x,(a,b))\mapsto(x,(a+s(x),b+s(x)))$, where $s(x)$ is
the sum of the coordinates of $x$, is regarded as a linear equivalence of the product of
the $m$-coordinate space with $\mathbb{F}_2^2$ to itself.
\end{definition}

\begin{theorem}[Action of the tail shear]
\lean{FABL.tailShearLinearEquiv_apply}
\label{FABL.tailShearLinearEquiv_apply}
\leanok
\uses{FABL.F₂Cube, FABL.tailShearLinearEquiv, FABL.𝔽₂}
\difficulty{1}
For every natural number $m$, every vector $x\\in(\\mathbb{F}_2)^m$, and all
$a,b\\in\\mathbb{F}_2$, the tail shear linear equivalence sends $(x,(a,b))$ to
$(x,(a+\\sum_{i=1}^m x_i,\\ b+\\sum_{i=1}^m x_i))$.
\end{theorem}

\begin{definition}[Paired coordinate equivalence]
\lean{FABL.pairedCoordinateEquiv}
\label{FABL.pairedCoordinateEquiv}
\leanok
\uses{FABL.F₂Cube, FABL.splitLastTwoLinearEquiv, FABL.tailShearLinearEquiv, FABL.𝔽₂}
For natural numbers $r$ and $k$, define a linear equivalence from the space of
$\mathbb{F}_2$-valued vectors with $r+2k$ coordinates to itself recursively. At $k=0$ it
is the identity; at $k+1$, split a vector into a block of $r+2k$ coordinates and its
final two coordinates, add the sum of the block’s coordinates to each final coordinate,
apply the equivalence for $k$ to the block, and reassemble the coordinates.
\end{definition}

\begin{theorem}[Effect on an appended coordinate pair]
\lean{FABL.pairedCoordinateEquiv_succ_appendTwo}
\label{FABL.pairedCoordinateEquiv_succ_appendTwo}
\leanok
\uses{FABL.F₂Cube, FABL.appendTwo, FABL.pairedCoordinateEquiv, FABL.splitLastTwoLinearEquiv, FABL.𝔽₂}
\difficulty{3}
Let $r,k$ be nonnegative integers, let $x\in\mathbb{F}_2^{r+2k}$, and let
$a,b\in\mathbb{F}_2$. The paired-coordinate linear automorphism on
$\mathbb{F}_2^{r+2(k+1)}$, applied to the vector formed by appending $a,b$ to $x$, gives
the vector formed by appending $a+\sum_i x_i$ and $b+\sum_i x_i$ to the image of $x$
under the paired-coordinate automorphism on $\mathbb{F}_2^{r+2k}$, where the sum is in
$\mathbb{F}_2$.
\end{theorem}

\begin{definition}[Left index in a coordinate pair]
\lean{FABL.pairedLeftIndex}
\label{FABL.pairedLeftIndex}
\leanok
For natural numbers $r$ and $k$ and an index $j\in\{0,\ldots,k-1\}$, the left index of
the $j$th pair among $r+2k$ coordinates is $r+2j$.
\end{definition}

\begin{definition}[Paired right index]
\lean{FABL.pairedRightIndex}
\label{FABL.pairedRightIndex}
\leanok
For natural numbers $r$ and $k$, and $j\in\{0,\ldots,k-1\}$, the paired right index is
the element $r+2j+1$ of $\{0,\ldots,r+2k-1\}$.
\end{definition}

\begin{theorem}[Compatibility of left indices under extension]
\lean{FABL.pairedLeftIndex_castSucc}
\label{FABL.pairedLeftIndex_castSucc}
\leanok
\uses{FABL.pairedLeftIndex}
\difficulty{1}
For nonnegative integers $r,k$ and an index $j\in\{0,\ldots,k-1\}$, the left
paired-coordinate index $r+2j$ among $k+1$ pairs, regarded as an index in the first
$r+2k$ coordinates of the larger index set, equals the left paired-coordinate index
$r+2j$ among $k$ pairs.
\end{theorem}

\begin{theorem}[Right index under successor embedding]
\lean{FABL.pairedRightIndex_castSucc}
\label{FABL.pairedRightIndex_castSucc}
\leanok
\uses{FABL.pairedRightIndex}
\difficulty{1}
Let $r,k$ be nonnegative integers and let $j\in\{0,\ldots,k-1\}$. The right index
associated with the successor embedding of $j$ into the next finite index set is the
image, under the natural inclusion that adds two coordinates, of the right index
associated with $j$; that is, $r+2(j+1)+1$ in the larger coordinate set corresponds to
$r+2j+1$ in the original set.
\end{theorem}

\begin{theorem}[Last left paired-coordinate index]
\lean{FABL.pairedLeftIndex_last}
\label{FABL.pairedLeftIndex_last}
\leanok
\uses{FABL.pairedLeftIndex}
\difficulty{3}
For nonnegative integers $r,k$, the left paired-coordinate index associated with the
final index among $k+1$ pairs is $r+2k$.
\end{theorem}

\begin{theorem}[Last right paired index]
\lean{FABL.pairedRightIndex_last}
\label{FABL.pairedRightIndex_last}
\leanok
\uses{FABL.pairedRightIndex}
\difficulty{3}
For nonnegative integers $r$ and $k$, the right index of the last pair among the $k+1$
pairs is the coordinate with offset $1$ in the final two-coordinate block after the
first $r+2k$ coordinates.
\end{theorem}

\begin{theorem}[Preservation of existing left coordinates]
\lean{FABL.pairedCoordinateEquiv_succ_left_castSucc}
\label{FABL.pairedCoordinateEquiv_succ_left_castSucc}
\leanok
\uses{FABL.F₂Cube, FABL.appendTwo, FABL.pairedCoordinateEquiv, FABL.pairedCoordinateEquiv_succ_appendTwo, FABL.pairedLeftIndex, FABL.𝔽₂}
\difficulty{3}
Let $r,k$ be nonnegative integers, $x\in\mathbb{F}_2^{r+2k}$, $a,b\in\mathbb{F}_2$, and
$j\in\{0,\ldots,k-1\}$. At the coordinate $r+2j$ in the transformed vector obtained by
appending $a,b$ to $x$ and applying the paired-coordinate automorphism for $k+1$ pairs,
the value equals the value at coordinate $r+2j$ of the image of $x$ under the
paired-coordinate automorphism for $k$ pairs.
\end{theorem}

\begin{theorem}[Preservation of existing right coordinates]
\lean{FABL.pairedCoordinateEquiv_succ_right_castSucc}
\label{FABL.pairedCoordinateEquiv_succ_right_castSucc}
\leanok
\uses{FABL.F₂Cube, FABL.appendTwo, FABL.pairedCoordinateEquiv, FABL.pairedCoordinateEquiv_succ_appendTwo, FABL.pairedRightIndex, FABL.𝔽₂}
\difficulty{3}
Let $r,k$ be nonnegative integers, $x\in\mathbb{F}_2^{r+2k}$, $a,b\in\mathbb{F}_2$, and
$j\in\{0,\ldots,k-1\}$. At the coordinate $r+2j+1$ in the transformed vector obtained by
appending $a,b$ to $x$ and applying the paired-coordinate automorphism for $k+1$ pairs,
the value equals the value at coordinate $r+2j+1$ of the image of $x$ under the
paired-coordinate automorphism for $k$ pairs.
\end{theorem}

\begin{theorem}[Value at the new left coordinate]
\lean{FABL.pairedCoordinateEquiv_succ_left_last}
\label{FABL.pairedCoordinateEquiv_succ_left_last}
\leanok
\uses{FABL.F₂Cube, FABL.appendTwo, FABL.pairedCoordinateEquiv, FABL.pairedCoordinateEquiv_succ_appendTwo, FABL.pairedLeftIndex, FABL.𝔽₂}
\difficulty{3}
Let $r,k$ be nonnegative integers, $x\in\mathbb{F}_2^{r+2k}$, and $a,b\in\mathbb{F}_2$.
The coordinate at index $r+2k$ of the image of the vector formed by appending $a,b$ to
$x$ under the paired-coordinate automorphism for $k+1$ pairs is $a+\sum_i x_i$, with the
sum taken in $\mathbb{F}_2$.
\end{theorem}

\begin{theorem}[Value at the new right coordinate]
\lean{FABL.pairedCoordinateEquiv_succ_right_last}
\label{FABL.pairedCoordinateEquiv_succ_right_last}
\leanok
\uses{FABL.F₂Cube, FABL.appendTwo, FABL.pairedCoordinateEquiv, FABL.pairedCoordinateEquiv_succ_appendTwo, FABL.pairedRightIndex, FABL.𝔽₂}
\difficulty{3}
Let $r,k$ be nonnegative integers, $x\in\mathbb{F}_2^{r+2k}$, and $a,b\in\mathbb{F}_2$.
The coordinate at index $r+2k+1$ of the image of the vector formed by appending $a,b$ to
$x$ under the paired-coordinate automorphism for $k+1$ pairs is $b+\sum_i x_i$, with the
sum taken in $\mathbb{F}_2$.
\end{theorem}

\begin{definition}[Paired index]
\lean{FABL.pairedIndex}
\label{FABL.pairedIndex}
\leanok
\uses{FABL.pairedLeftIndex, FABL.pairedRightIndex}
For natural numbers $r$ and $k$, define an index on the disjoint union of two copies of
$\{0,\ldots,k-1\}$. It sends the first copy's index $j$ to $r+2j$ and the second copy's
index $j$ to $r+2j+1$, in each case as an element of $\{0,\ldots,r+2k-1\}$.
\end{definition}

\begin{theorem}[Injectivity of paired coordinate indices]
\lean{FABL.pairedIndex_injective}
\label{FABL.pairedIndex_injective}
\leanok
\uses{FABL.pairedIndex, FABL.pairedLeftIndex, FABL.pairedRightIndex}
\difficulty{5}
For nonnegative integers $r,k$, the map from the disjoint union of two copies of
$\{0,\ldots,k-1\}$ to $\{0,\ldots,r+2k-1\}$ that sends the first copy's index $j$ to
$r+2j$ and the second copy's index $j$ to $r+2j+1$ is injective.
\end{theorem}

\begin{definition}[Projection onto the prefix coordinates]
\lean{FABL.prefixLastTwoLinearMap}
\label{FABL.prefixLastTwoLinearMap}
\leanok
\uses{FABL.F₂Cube, FABL.splitLastTwoLinearEquiv, FABL.𝔽₂}
For a natural number $m$, the linear map from $(\mathbb{F}_2)^{m+2}$ to
$(\mathbb{F}_2)^m$ sends a vector to its first $m$ coordinates, discarding its final two
coordinates.
\end{definition}

\begin{theorem}[Projection after appending two coordinates]
\lean{FABL.prefixLastTwoLinearMap_appendTwo}
\label{FABL.prefixLastTwoLinearMap_appendTwo}
\leanok
\uses{FABL.F₂Cube, FABL.appendTwo, FABL.prefixLastTwoLinearMap, FABL.𝔽₂}
\difficulty{1}
Let $m$ be a natural number, let $x$ be a vector in $\mathbb{F}_2^m$, and let
$a,b\in\mathbb{F}_2$. The linear projection from $\mathbb{F}_2^{m+2}$ onto its first $m$
coordinates sends the vector formed by appending $a,b$ to $x$.
\end{theorem}

\begin{definition}[Recursive paired linear form]
\lean{FABL.pairedAffineLinearForm}
\label{FABL.pairedAffineLinearForm}
\leanok
\uses{FABL.F₂Cube, FABL.coordinateSumLinearMap, FABL.prefixLastTwoLinearMap, FABL.𝔽₂}
Let $r$ be a natural number, and let $L_{r,k}$ be a linear functional on
$(\mathbb{F}_2)^{r+2k}$. Set $L_{r,0}=0$. Recursively, for
$x\in(\mathbb{F}_2)^{r+2k+2}$, set $L_{r,k+1}(x)=L_{r,k}(x_{\le
r+2k})+\sum_{i=0}^{r+2k-1}(x_{\le r+2k})_i$, where $x_{\le r+2k}$ denotes the vector
formed by the first $r+2k$ coordinates of $x$ and the sum is in $\mathbb{F}_2$.
\end{definition}

\begin{theorem}[Paired affine form after appending two coordinates]
\lean{FABL.pairedAffineLinearForm_succ_appendTwo}
\label{FABL.pairedAffineLinearForm_succ_appendTwo}
\leanok
\uses{FABL.F₂Cube, FABL.appendTwo, FABL.pairedAffineLinearForm, FABL.𝔽₂}
\difficulty{1}
Let $r,k$ be natural numbers, let $x$ be a vector of length $r+2k$ over the field with
two elements, and let $a,b$ lie in that field. The paired affine linear form with
parameter $k+1$, evaluated on the vector formed by appending $a,b$ to $x$, equals the
paired affine linear form with parameter $k$ evaluated on $x$, plus the sum of the
coordinates of $x$.
\end{theorem}

\begin{theorem}[Full-sum formula for the quadratic bit]
\lean{FABL.completeQuadraticBit_eq_full_sum}
\proofsource{Analysis-of-Boolean-Functions-by-Ryan-ODonnell}{pdf-ffbe9c7b3306-p034-b009}
\label{FABL.completeQuadraticBit_eq_full_sum}
\leanok
\uses{FABL.F₂Cube, FABL.completeQuadraticBit}
\difficulty{4}
Let $m$ be a nonnegative integer and let $x\in\mathbb{F}_2^m$. The complete quadratic
bit of $x$ equals $\sum_{i=0}^{m-1}\sum_{j=0}^{m-1} [i<j]x_i x_j$, where each summand is
$x_i x_j$ when $i<j$ and $0$ otherwise, and the sum is in $\mathbb{F}_2$.
\end{theorem}

\begin{theorem}[Quadratic bit of an appended vector]
\lean{FABL.completeQuadraticBit_append}
\label{FABL.completeQuadraticBit_append}
\leanok
\uses{FABL.F₂Cube, FABL.completeQuadraticBit, FABL.completeQuadraticBit_eq_full_sum}
\difficulty{5}
Let $m,p$ be nonnegative integers, and let $x\in\mathbb{F}_2^m$ and
$y\in\mathbb{F}_2^p$. The complete quadratic bit of the vector formed by appending $y$
to $x$ equals the complete quadratic bit of $x$, plus that of $y$, plus $(\sum_i
x_i)(\sum_j y_j)$, with all operations in $\mathbb{F}_2$.
\end{theorem}

\begin{theorem}[Quadratic bit in two coordinates]
\lean{FABL.completeQuadraticBit_two}
\label{FABL.completeQuadraticBit_two}
\leanok
\uses{FABL.completeQuadraticBit, FABL.completeQuadraticBit_eq_full_sum, FABL.𝔽₂}
\difficulty{3}
For all $a,b\in\mathbb{F}_2$, the complete quadratic bit of the two-coordinate vector
$(a,b)$ is $ab$.
\end{theorem}

\begin{theorem}[Complete quadratic bit append two]
\lean{FABL.completeQuadraticBit_appendTwo}
\statementsource{Analysis-of-Boolean-Functions-by-Ryan-ODonnell}{pdf-ffbe9c7b3306-p166-b002}
\label{FABL.completeQuadraticBit_appendTwo}
\leanok
\uses{FABL.F₂Cube, FABL.appendTwo, FABL.completeQuadraticBit, FABL.completeQuadraticBit_append, FABL.completeQuadraticBit_two, FABL.𝔽₂}
\difficulty{3}
% dataset-pipeline: stub -- prose pending build_informal_statements.py
Statement pending regeneration.
\end{theorem}

\begin{theorem}[Paired product sum after appending two coordinates]
\lean{FABL.pairedProductSum_succ_appendTwo}
\label{FABL.pairedProductSum_succ_appendTwo}
\leanok
\uses{FABL.F₂Cube, FABL.appendTwo, FABL.pairedCoordinateEquiv, FABL.pairedCoordinateEquiv_succ_left_castSucc, FABL.pairedCoordinateEquiv_succ_left_last, FABL.pairedCoordinateEquiv_succ_right_castSucc, FABL.pairedCoordinateEquiv_succ_right_last, FABL.pairedLeftIndex, FABL.pairedRightIndex, FABL.𝔽₂}
\difficulty{1}
Let $r,k$ be nonnegative integers, let $x\in\mathbb F_2^{r+2k}$, and let $a,b\in\mathbb
F_2$. Append $a,b$ to $x$ to obtain a vector in $\mathbb F_2^{r+2(k+1)}$, and apply the
paired-coordinate linear equivalence at level $k+1$. The sum of the products of the
transformed coordinates in each of the $k+1$ designated pairs equals the corresponding
sum at level $k$ for $x$, plus
$\bigl(a+\sum_{i=1}^{r+2k}x_i\bigr)\bigl(b+\sum_{i=1}^{r+2k}x_i\bigr)$, with all
arithmetic in $\mathbb F_2$.
\end{theorem}

\begin{theorem}[Paired decomposition of the quadratic bit]
\lean{FABL.completeQuadraticBit_paired_decomposition}
\proofsource{Analysis-of-Boolean-Functions-by-Ryan-ODonnell}{pdf-ffbe9c7b3306-p166-b002}
\label{FABL.completeQuadraticBit_paired_decomposition}
\leanok
\uses{FABL.F₂Cube, FABL.appendTwo, FABL.completeQuadraticBit, FABL.completeQuadraticBit_appendTwo, FABL.pairedAffineLinearForm, FABL.pairedAffineLinearForm_succ_appendTwo, FABL.pairedCoordinateEquiv, FABL.pairedLeftIndex, FABL.pairedProductSum_succ_appendTwo, FABL.pairedRightIndex, FABL.splitLastTwoLinearEquiv, FABL.𝔽₂}
\difficulty{6}
Let $r$ and $k$ be nonnegative integers with $r\le 1$, and let
$x\in\mathbb{F}_2^{r+2k}$. Define linear coordinate changes $T_{r,k}$ recursively:
$T_{r,0}$ is the identity; to form $T_{r,k+1}$, split off the last two coordinates, add
the sum of the preceding coordinates to each of those two, apply $T_{r,k}$ to the
preceding coordinates, and reassemble. Define linear forms $L_{r,0}=0$ and
$L_{r,k+1}(x)=L_{r,k}(x_{<r+2k})+\sum_{i<r+2k}x_i$, where $x_{<r+2k}$ is the prefix of
length $r+2k$. Then, in $\mathbb{F}_2$, $\sum_{i<j}x_i
x_j=L_{r,k}(x)+\sum_{j=0}^{k-1}(T_{r,k}x)_{r+2j}(T_{r,k}x)_{r+2j+1}$.
\end{theorem}

\begin{theorem}[Quadratic form as a sum of paired products]
\lean{FABL.exists_completeQuadraticBit_pairing_of_remainder}
\proofsource{Analysis-of-Boolean-Functions-by-Ryan-ODonnell}{pdf-ffbe9c7b3306-p166-b002}
\label{FABL.exists_completeQuadraticBit_pairing_of_remainder}
\leanok
\uses{FABL.F₂Cube, FABL.IsF₂Linear, FABL.affineFunction, FABL.completeQuadraticBit, FABL.completeQuadraticBit_paired_decomposition, FABL.isF₂Linear_iff_exists_dotProduct, FABL.pairedAffineLinearForm, FABL.pairedCoordinateEquiv, FABL.pairedIndex, FABL.pairedIndex_injective, FABL.pairedLeftIndex, FABL.pairedRightIndex, FABL.𝔽₂}
\difficulty{5}
Let $r,k$ be natural numbers with $r\le1$. On the binary cube of dimension $r+2k$, the
quadratic function $Q(x)=\sum_{i<j}x_i x_j$ can be written as an affine function plus
$\sum_{j=1}^k L_j(x)R_j(x)$, where the $2k$ linear forms $L_1,R_1,\ldots,L_k,R_k$ are
linearly independent over $\mathbb F_2$.
\end{theorem}

\begin{theorem}[Transporting a quadratic decomposition]
\lean{FABL.completeQuadraticBit_decomposition_cast}
\label{FABL.completeQuadraticBit_decomposition_cast}
\leanok
\uses{FABL.F₂Cube, FABL.affineFunction, FABL.completeQuadraticBit, FABL.𝔽₂}
\difficulty{1}
Let $m,n,k$ be nonnegative integers with $m=n$. Suppose there are $b\in\mathbb{F}_2$, a
vector $a\in\mathbb{F}_2^m$, and two families of linear functionals $L_1,\ldots,L_k$ and
$R_1,\ldots,R_k$ on $\mathbb{F}_2^m$ such that the combined family is linearly
independent and, for every $x\in\mathbb{F}_2^m$, the complete quadratic bit equals
$b+a\cdot x+\sum_{j=1}^k L_j(x)R_j(x)$. Then there are such a scalar, vector, and two
families of linear functionals on $\mathbb{F}_2^n$, again with the combined family
linearly independent, giving this identity for every $x\in\mathbb{F}_2^n$.
\end{theorem}

\begin{theorem}[Exists complete quadratic bit affine independent decomposition]
\lean{FABL.exists_completeQuadraticBit_affine_independent_decomposition}
\statementsource{Analysis-of-Boolean-Functions-by-Ryan-ODonnell}{pdf-ffbe9c7b3306-p165-b009}
\label{FABL.exists_completeQuadraticBit_affine_independent_decomposition}
\leanok
\uses{FABL.F₂Cube, FABL.affineFunction, FABL.completeQuadraticBit, FABL.completeQuadraticBit_decomposition_cast, FABL.exists_completeQuadraticBit_pairing_of_remainder, FABL.𝔽₂}
\difficulty{4}
% dataset-pipeline: stub -- prose pending build_informal_statements.py
Statement pending regeneration.
\end{theorem}
